There are three files for every model and every forecast origin. \textbf{They are not three models
and not three runs.} There was one run per model. It produced 16 draws for every cell and every
month, and a point forecast needs one number rather than 16, so something has to collapse them.
There is more than one defensible way to do that, and three were used.

\begin{table}[h!]
\centering
\small
\renewcommand{\arraystretch}{1.3}
\caption{The three frames. One run, one set of draws, collapsed three ways.}
\label{tab:frames}
\begin{tabular}{@{}p{1.1cm}p{3.6cm}p{9.3cm}@{}}
\toprule
\textbf{Frame} & \textbf{Directory} & \textbf{What it is} \\ \midrule
I   & \texttt{I\_mean} &
  The arithmetic mean of the 16 draws, taken in double precision on the count scale.
  \textbf{This is the estimator the models declare}, and the only one of the three that the
  delivery pipeline produces. \\ \addlinespace
II  & \texttt{II\_q95} &
  The 95th percentile across the 16 draws. An upper estimate: the magnitude that is exceeded in
  roughly one draw in twenty. \\ \addlinespace
III & \texttt{III\_conditional\_mean} &
  $\mathrm{E}[y \mid y > 0]$ --- the mean taken over the positive draws only, equivalently
  $\mathrm{E}[y] / P(y>0)$. The expected magnitude \emph{given} that something happens, rather
  than the expected magnitude including the chance that nothing does. \\
\bottomrule
\end{tabular}
\end{table}

\noindent Lexical sorting puts these in the order III, II, I, so a file browser or a Drive listing
will show them backwards. The Roman numerals are in the directory names so that the intended
reading order survives being sorted.

\subsection*{Why the Collapse Happens Outside the Run}

All eight models declare \texttt{evaluation\_mode: stochastic}, which tells the run to keep the
full set of draws and not to reduce them. Each model also declares an aggregation method
(\texttt{arithmetic\_mean}), but under stochastic evaluation that declaration has no effect and the
run says so in its own log. The consequence is that the collapse is a separate, later step, taken
deliberately rather than as a side effect of how the run was configured --- and that is the step
that produced these three frames.

\subsection*{What Frame I Can Be Checked Against, and What Frames II and III Cannot}

This needs stating plainly, because it is the one asymmetry in the delivery.

\textbf{Frame I is reproducible.} It is produced by a collapse step held in views-models, at
\texttt{tools/collapse/}\allowbreak\texttt{collapse\_predictions.py}, which implements the
arithmetic mean and the median and refuses anything else. It takes the mean in double precision, on raw
counts, and refuses the input outright on a non-finite value, a negative value, a disagreement in
draw count between targets, a duplicated cell-month, or a maximum low enough to suggest the values
are still on a log scale. Those refusals are tested.

\textbf{Frames II and III are not reproducible, and no code for them exists.} Neither the
percentile nor the conditional mean is implemented anywhere in views-models; both were computed
once, in an ad-hoc step, and the posterior draws they were computed from are not on any machine we
still have. The draws were written and compressed on the rented GPUs --- the run manifests record
13~GB of raw draws per model compressing to roughly 9~MB --- but they were never transferred off,
and the machines are gone. \textbf{So frames II and III cannot be recomputed, re-derived or
checked against the posterior.} They are what they are: a probe, recorded here as such.

This also corrects a claim made in my own cost memo of 28 September, which stated that the full
posterior had been retained and that a different summary statistic could therefore be produced
without a new run. That is not true on disk. \textbf{Any statistic other than these three needs a
new GPU run.}

\subsection*{Two Properties of These Estimators}

\textbf{Frame II is coarse at 16 draws.} The 95th percentile of 16 values falls between the 15th
and 16th order statistic, so a single draw moves it, and at a sample this small an upper quantile
is biased low. It is reported as computed, with no small-sample correction. If the draw count is
ever raised, frame II must be re-measured before anything relies on it.

\textbf{Frame III inflates the sparse targets most.} Dividing by $P(y>0)$ multiplies a cell that
was positive in one draw of 16 by about 16. The effect is largest where positive draws are rarest,
which is why frame III runs 4.2 times frame I on state-based fatalities but 7.5 times on one-sided
(Section~\ref{sec:magnitude}). One ratio quoted for all three targets would hide that.

\subsection*{How the Three Are Meant to Be Used}

All three go into the selection routine as if they were separate models, and the published
selection criteria decide. That is settled. I am not proposing one of them here, and the frames
are not ranked anywhere in this note.

One provenance correction while the frames are in view: the README uploaded with the folder cites
``ADR-023 Amendment 1'' as the authority for the three frames. No such document exists. ADR-023
governs the collapse and was accepted on 28 September 2026; its only amendment is dated 6 October
2026 and concerns an unrelated library. The authority for frame I is ADR-023 and the commit on the
title page. Frames II and III have no governing document, which is part of why this section says
what it says about them.
